\subsection{Caso 3. Exposición de servicios HTTP, FTP y otros servicios heredados}
\label{caso:servicios-heredados}

\begin{situacion}
El servidor presenta servicios abiertos que podrían estar innecesarios o débilmente configurados. Los riesgos de aplicaciones web se pueden contrastar con el catálogo OWASP \parencite{owasp2021}.
\end{situacion}

\paragraph{Procedimiento.}
\begin{itemize}[leftmargin=*, topsep=2pt, itemsep=2pt]
	\item Inventariar servicios y puertos en el servidor propio usando el inventario del sistema y una comprobación autorizada de alcance limitado.
	\item Revisar los registros de autenticación y los eventos del sistema; documentar versión, necesidad operativa y propietario de cada servicio.
	\item Comprobar si se utiliza cifrado y si las cuentas siguen el principio de mínimo privilegio.
\end{itemize}

\paragraph{Indicadores que deben buscarse.}
Servicios expuestos sin necesidad, software sin soporte, intentos de autenticación fallidos, uso de protocolos en texto claro o configuraciones por defecto.

\paragraph{Mitigación y defensa.}
Deshabilitar servicios no requeridos; preferir SFTP/SSH o HTTPS frente a protocolos en texto claro; aplicar MFA donde corresponda; parchear y usar cuentas de mínimo privilegio.

\paragraph{Evidencias que se entregan.}
Inventario de servicios, eventos relevantes, matriz riesgo/criticidad y plan de remediación.

\paragraph{Preguntas de análisis.}
\begin{enumerate}[leftmargin=*, topsep=2pt, itemsep=2pt]
	\item ¿Cómo se prioriza un puerto abierto?
	\item ¿Qué diferencia hay entre exposición, vulnerabilidad y explotación?
\end{enumerate}

\begin{table}[H]
	\caption{Inventario de servicios expuestos}
	\label{tab:inventario-servicios}
	\centering
	\footnotesize
	\renewcommand{\arraystretch}{1.25}
	\begin{tabularx}{\textwidth}{@{}c >{\raggedright\arraybackslash}p{0.15\textwidth} >{\raggedright\arraybackslash}p{0.12\textwidth} c >{\raggedright\arraybackslash}p{0.18\textwidth} Y@{}}
		\toprule
		\textbf{Puerto} & \textbf{Servicio} & \textbf{Versión} & \textbf{¿Necesario?} & \textbf{Propietario} & \textbf{Riesgo / observación} \\
		\midrule
		& & & & & \\
		& & & & & \\
		& & & & & \\
		& & & & & \\
		& & & & & \\
		& & & & & \\
		\bottomrule
	\end{tabularx}
\end{table}

\registrohallazgo{Caso 3}
